\documentclass[12pt,a4paper]{article}
\usepackage{geometry}
\geometry{left=2.5cm,right=2.5cm,top=2.0cm,bottom=2.5cm}
\usepackage[english]{babel}
\usepackage{amsmath,amsthm}
\usepackage{amsfonts}
\usepackage[longend,ruled,linesnumbered]{algorithm2e}
\usepackage{fancyhdr}
\usepackage{ctex}
\usepackage{array}
\usepackage{listings}
\usepackage{color}
\usepackage{graphicx}

\begin{document}


\title{
{《优化理论与算法》第 5 次作业
}
}
\date{}

% \author{
% 姓名：\underline{你的名字}~~~~~~
% 学号：\underline{你的学号}~~~~~~}

\maketitle

\begin{itemize}
	\item 作业截止于{\bf 周五（4/10/2026）}，在Canvas系统中提交PDF或照片文件
	\item 作业题目的tex文件可以在Canvas中下载
	\item 鼓励相互讨论，但不要抄袭.
\end{itemize}
%%%%%%%%%%%%%%%%%%%%%%%%%%%%%%%%%%%%%%%%%%%%%%%%%%%%%%%%%%%%%%%%%
%%%%%%%%%%%%%%%%%%%%%%%%%%%%%%%%%%%%%%%%%%%%%%%%%%%%%%%%%%%%%%%%%

\section{阅读与积累}
\paragraph{1.1}
\begin{itemize}
	\item 阅读Bertsimas and Tsitsikils，也就是Introduction to Linear Optimization教材，第4.1-4.4节、5.1-5.2节相关内容。
\end{itemize}

\section{习题}

\paragraph{2.1}
考虑以下线性规划问题：
\begin{align*}
\min \;&\quad x_1 - x_2 \\
\text{s.t.} \;&\quad2x_1 + 3x_2 - x_3 + x_4 \leq 0 \\
 \;&\quad3x_1 + x_2 + 4x_3 - 2x_4 \geq 3 \\
 \;&\quad-x_1 - x_2 + 2x_3 + x_4 = 6 \\
 \;&\quad x_1 \leq 0, x_2 \text{ 无约束 }, x_3 \geq 0, x_4 \text{ 无约束 }
\end{align*}

写出相应的对偶问题。

\paragraph{2.2}
考虑以下线性规划问题：

\begin{align*}
   \max \quad &\; 3x_1 + 2x_2\\
   s.t.\quad &\;x_1 + x_2 \leq 4 \\
&\;2x_1 + x_2 \leq 5 \\
&\;x_1, x_2 \geq 0
\end{align*}
利用互补松弛条件判断下列解是否是最优解
\begin{itemize}
    \item 判断解\( x_1 = 1 \)，\( x_2 = 3 \)是否是最优解
    \item 判断解\( x_1 = 2 \)，\( x_2 = 1 \)是否是最优解
\end{itemize}

\paragraph{2.3}
试证明：线性系统$Ax=b, x \geq 0$无解，只需找的某可行解$y$使之满足$A^Ty\geq0, b^Ty<0$.


\paragraph{2.4}
(1)利用单纯形法求解如下线性规划
\begin{equation*}
\begin{array}{lllll@{}ll}
\text{maximize}  & 2x_1+4x_2+x_3+x_4 &\\
\text{subject to}& x_1+3x_2 +x_4    &\leq 4 \\
&2x_1 +x_2   &\leq 3  \\
&x_2 +4x_3+x_4   &\leq 3  \\
&x_1,x_2,x_3,x_4 \geq 0 &  
\end{array}
\end{equation*}

(2)在（1）求解的最优基的基础上，$b=[4, 3, 3]^T$ 中每个元素在什么范围内变动，可以保持最优基不变呢？（假设某一个元素变动时，其他元素不变）

(3)在（1）求解的最优解的基础上，$c=[2,4,1,1]^T$中每个元素在什么范围内变动，可以保持最优解不变呢？（假设某一个元素变动时，其他元素不变）



 \paragraph{2.5}
考虑课件dual3对偶理论的应用，P26展示了排产问题(P27求解到最优解的单纯形表)
 \begin{enumerate}
 	\item 若新增某生产产品，使得A矩阵新增一列，并带来新的利润。此时，原问题改变为
 	\begin{equation*}
 	\begin{array}{lllll@{}ll}
 	\text{maximize}  & x_1&+2x_2 &+3x_3&\\
\text{subject to}& x_1 &&+x_3    &\leq 100 \\
	&  &2x_2  &+x_3  &\leq 200 \\
 	&x_1 &+x_2  &+x_3  &\leq 150  \\
 	&x_1 \geq 0, &x_2 \geq 0, &x_3 \geq 0& 
 	\end{array}
 	\end{equation*}
 	请利用灵敏度分析，找到如上问题最优解。
	
	
 	\item 若新增产能约束，使得A矩阵新增一行。此时，原问题改变为
 	\begin{equation*}
 	\begin{array}{lllll@{}ll}
 	\text{maximize}  & x_1&+2x_2 &\\
 	\text{subject to}& x_1 &  &\leq 100 \\
 	&  &2x_2    &\leq 200 \\
	&x_1 &+x_2   &\leq 150  \\
 	&4x_1 &+x_2   &\leq 250  \\
	&x_1 \geq 0, &x_2 \geq 0 &
 	\end{array}
 	\end{equation*}
	请利用灵敏度分析，找到如上问题最优解。
		
	
 \end{enumerate}




\end{document}